\section{\acs{NEF}}
\label{chap:proposal:nef}

As outlined in Section \ref{sec:sba:nef}, the \acf{NEF} is the main \acs{5G} core network function responsible for exposing network capabilities to external (and possibly internal) consumers. To accomplish this objective, it offers a set of standardized northbound \acsp{API} that mediate access between the \acs{5G} core and consuming \acsp{AF}, while abstracting internal network functions and their implementation details.

In the context of the proposed testbed architecture, three \acs{NEF} northbound \acsp{API} are considered to address distinct, yet complementary aspects of network application behavior. The \textit{AsSessionWithQoS} \acs{API} targets \acs{QoS} control, providing applications with the means to establish and manage application sessions with specific \acs{QoS} profiles designed to guarantee, for example, specific bandwidth or latency requirements~\cite{ETSI_TS_129_122_V18_8_0}. However, passively relying on the \acs{QoS} session established via this \acs{API} is not sufficient, as undetected violations of network requirements directly risk the correct functioning of the Network Application. Therefore, there is a clear need for a mechanism capable of actively monitoring the network and detecting any violations of the defined \acs{QoS} in real time. The \textit{AnalyticsExposure} \acs{API} addresses this problem giving developers the ability to query network analytics and subscribe to analytics notifications~\cite{ETSI_TS_129_522_V18_11_0}, allowing applications to monitor network metrics and adapt their behavior according to the observed conditions.

\acs{QoS} violations and changes in network conditions are not the only criteria that can require adaptive application behavior. According to its use case, the network application may need to be notified about specific events in the \acs{3GPP} network, completely unrelated to \acs{QoS}, such as an \acs{UE} becoming unreachable or changing location. The \textit{MonitoringEvent} \acs{API} was designed to create and manage such monitoring event subscriptions~\cite{ETSI_TS_129_122_V18_8_0}. The following subsections describe each \acs{API} and its corresponding procedure, detailing the interactions with the necessary \acsp{NF} to achieve the specified use cases.\review{make this valid}


\subsection{AsSessionWithQoS \acs{API}}

The \textit{AsSessionWithQoS} \acs{API}, defined in \acs{3GPP} TS~29.122~\cite{ETSI_TS_129_122_V18_8_0}, allows applications to configure sessions indicating the desired level of \acs{QoS} for a given traffic flow. However, passively relying on the \acs{QoS} session established via this \acs{API} is not sufficient, as undetected violations of \acs{QoS} requirements directly risk the correct functioning of the Network Application. Therefore, in order to ensure better service experience and avoid service interruption, the \textit{AsSessionWithQoS} \acs{API} also provides multiple mechanisms to detect and help mitigate the effects of \acs{QoS} downgrades and network anomalies.

The interaction is initiated when the \acs{AS}/\acs{AF} (e.g. a network application) submits an HTTP POST request to the \acs{NEF}, carrying an \texttt{AsSessionWithQoSSubscription} resource body~\cite{ETSI_TS_129_122_V18_8_0}. The subscription body is composed by an extensive set of parameters, including the target \acs{UE}(s) and traffic flow, the desired \acs{QoS} and monitoring and notification settings. As it is typical of \acs{3GPP} resource descriptions, these parameters are governed by conditional notes that dictate whether a parameter shall be included according to the rest of the payload parameters and the supported \acs{API} features (some fields can only be supplied if the \acs{5GC} supports a given feature). The purpose of this section is not to provide an exhaustive analysis of the spec and all its parameters and constraints~\cite{ETSI_TS_129_122_V18_8_0}, but instead to highlight how the northbound \acs{API} payload can be constructed to express the desired \acs{QoS} configuration, with a particular focus on the following key attributes: 
\begin{itemize}
    \item Notification Destination: The parameter \texttt{notificationDestination} contains the \acs{URL} to receive notification events. This field corresponds to the only mandatory parameter in this payload with the remaining optional parameters being constrained by restrictions defined in the \texttt{AsSessionWithQoSSubscription} data type definition.
    \item Target \acs{UE}(s): When the "GMEC" feature is not supported, either (i) a single target \acs{UE} is specified via exactly one of its \acs{IP} address (\texttt{ueIpv4Addr} or \texttt{ueIpv6Addr}) or \acs{MAC} address (\texttt{macAddr}); or (ii) multiple target \acsp{UE} are specified with \texttt{listUeAddrs} parameter. When "GMEC" feature is supported and the target \acs{UE}(s) are not identified by \acs{UE} addresses (via the previous parameters), either the \texttt{gpsi} or the \texttt{extGroupId} attribute can be provided, identifying an unique \acs{UE} or group of \acsp{UE}, respectively;
    \item Traffic Flow: The traffic flow subject to \acs{QoS} treatment is specified differently depending on the target \acs{UE} identification method described previously. It corresponds to an IP flow defined through \texttt{flowInfo} if an IP (IPV4 or IPV6) address was provided or to an Ethernet flow described using \texttt{ethFlowInfo} (or \texttt{enEthFlowInfo} if the feature "EnEthAsSessionQoS\_5G" is supported) if a \acs{MAC} address was supplied and the "AppId" feature is not supported;
    \item \acs{QoS} Configuration: Having identified the target \acs{UE}(s) and the associated traffic flow, the payload must then specify what \acs{QoS} should be enforced on that path. Generally, this can be achieved in one of the following ways, depending on the features supported by the \acs{API}: (i) by supplying references to pre-defined \acs{QoS} information (e.g. \acs{QCI}/\acs{5QI} and associated parameters defined by the operator) through the \texttt{qosReference} attribute for the primary \acs{QoS} configuration, optionally complemented by \texttt{altQoSReferences}, for an ordered list of alternatives ranked in ascending order of priority; (ii) by specifying \texttt{altQosReqs} detailing a list of alternative service requirements that include individual \acs{QoS} parameter sets for guaranteed bandwidth in uplink (\texttt{gbrUl}), guaranteed bandwidth in downlink (\texttt{gbrD}), the packet delay budget (\texttt{pdb}) and the packet error rate (\texttt{per}). \texttt{tscQosReq} field may also be used to specify \acs{QoS} requirements specific to time sensitive communications;
    \item \acs{QoS} Signalling: The \texttt{disUeNotif} attribute controls whether to disable signalling of \acs{QoS} flow parameters to the \acs{UE}, when its \acs{QoS} conditions change to either the \acs{QoS} profile or an Alternative \acs{QoS} Profile;
    \item Usage Threshold: The field \texttt{usageThreshold} defines the time period (\texttt{duration}) and/or traffic volume (\texttt{totalVolume}, \texttt{downlinkVolume} or \texttt{uplinkVolume}) in which the \acs{QoS} is to be applied;
    \item \acs{QoS} Monitoring: Depending on the subset of monitoring features supported by the \acs{API} ("QoSMonitoring\_5G", "EnQoSMonGMEC", "EnQoSMonListUE\_5GGMEC") the following monitoring information can be requested: packet delay (\texttt{qosMonInfo}), packet delay variation (\texttt{pdvMon}), data rate (\texttt{qosMonDatRate}) and congestion measurements (\texttt{qosMonConReq}). Each set of metrics can be reported either periodically or based on events, which are triggered when the requested thresholds are crossed.
\end{itemize}

Following the creation request, the \acs{NEF} maps the supplied \acs{QoS} parameters and coordinates with the \acs{PCF} to enforce the policy.

responds to the HTTP POST with the created \texttt{AsSessionWithQoSSubscription} resource, returning its \acs{URI} for subsequent management operations in the \texttt{self} parameter. 


\plan{
    \begin{enumerate}
        \item Define the purpose of the API. Do not be repetitive comparetively to the first part of this section. At least rephrase it in a different way that doens't saturate the reader.
        \item Describe the typical interaction flow: an application submits a QoS subscription request, provide the exact subscription request type (identifying the UE, IP flow, and desired QoS profile - map these to the specific fields); the NEF maps this to internal 5G QoS parameters and coordinates with the PCF to enforce the policy.
        \item Highlight the key parameters exposed by the API (QoS reference identifiers, UE identifiers, flow descriptors, and notification endpoints) and clarify how they abstract away the complexity of internal QoS class mappings.
        \item Position this API as the primary mechanism enabling network applications to express and enforce latency/bandwidth requirements, directly supporting use-case-specific network conditions.
    \end{enumerate}
}

\subsection{MonitoringEvent \acs{API}}

\plan{
    \begin{enumerate}
        \item Define the purpose of the API
        \item Enumerate the main event types relevant to the work: location reporting, loss of connectivity, UE reachability and roaming status
        \item Describe the typical interaction flow: applications register a monitoring event configuration (specifying the target UE(s), event type, reporting frequency, and callback URI); the NEF then triggers notifications to the application upon event occurrence.
        \item Frame MonitoringEvent \acs{API} as a complementary mechanism to the previously mentioned AsSessionWithQoS \acs{API}, as it enables network applications to detect conditions (based on \acs{UE} state) that warrant \acs{QoS} changes through the AsSessionWithQoS \acs{API}.
    \end{enumerate}
}

\subsection{AnalyticsExposure \acs{API}}

\plan{
    \begin{enumerate}
        \item Define the purpose of the API 
        \item Highlight the relevance of the \acs{API} regarding network observability and the range of metrics mentioned in the next section.
        \item Emphasize the role of the \acs{NWDAF} as the native \acs{5G} mechanism for network observability, being responsible for gathering metrics scattered across multiple network functions, and deriving new ones.
        \item Reintroduce the \acs{NEF} analytics exposure \acs{API} as the standardized interface for exposing network analytics to external consumers, whether sourced from the \acs{NWDAF} or obtained via alternative means (e.g., directly from raw IP traffic analysis or via active probing). This is specially relevant, as it allows the validation of network application requirements and behavior without the need for an \acs{NWDAF}.
        \item Introduce some of the metrics that can be exposed via this \acs{API} (mention that a more detailed version of this will be included in the next chapter)
        \item Position this \acs{API} as the bridge between internal network observability and the network application layer, allowing both: the change of \acs{QoS} profile according to network conditions; 2. assurance of network's expected performance, as per the testbed client slice requirements; 3. lastly, the validation of application behavior and requirements
    \end{enumerate}
}